\documentclass[a5paper,10pt]{article}

% https://github.com/InDevRus/filippov-solutions
% Нумерация страниц
\pagenumbering{gobble}

% Настройка полей
\usepackage{geometry}
\geometry{left=1cm}
\geometry{right=1cm}
\geometry{top=1cm}
\geometry{bottom=1.5cm}

% Вставка таблицы
\usepackage{array}
\usepackage{tabularx}

% Инструменты выравнивания формул
\usepackage{amsmath}
\usepackage{mathtools}

% Диакритика в математике
\usepackage{amsfonts}

% Вычисления размеров на ходу
\usepackage{calc}

% Удобные записи производных
\usepackage[italic=true]{derivative}

% Настройки сносок
\usepackage{enotez}

% Длинные знаки векторов
\usepackage{anyfontsize}
\usepackage[b]{esvect}

% Вставка картинок
\usepackage{graphicx}

% Вставка красной строки
\usepackage{indentfirst}
\setlength{\parindent}{1em}

% Условия в командах
\usepackage{xifthen}

% Луа-скрипты
\usepackage{luacode}

% Настройка команд отдельно в математическом и текстовом режимах
\usepackage{mathcommand}

% Для повторения LaTeX кода
\usepackage{multido}

% Конфигурация отступов формул
\delimitershortfall-1sp
\usepackage{mleftright}
\mleftright

% Растягивание объектов
\usepackage{relsize}

% Обтекание текстом
\usepackage{wrapfig}
\usepackage{subcaption}
\usepackage{floatflt}

% Поддержка ввода кириллицы
\usepackage[english, russian]{babel}

% Настройка корневой позиции для компилятора
\usepackage{import}
\providecommand*{\ProjectRootPath}{..}

\import{\ProjectRootPath/preambles/macros/}{theorems}

\import{\ProjectRootPath/preambles/fonts}{font_setup}

% Знак градуса
\usepackage{siunitx}

\import{\ProjectRootPath/preambles/macros/}{operators}
\import{\ProjectRootPath/preambles/macros/}{redefinitions}

\import{\ProjectRootPath/preambles/fonts}{letters_setup}
\import{\ProjectRootPath/preambles/fonts/adjustments}{custom_superscripts}

\import{\ProjectRootPath/preambles/custom}{formatting}
\import{\ProjectRootPath/preambles/custom}{frames}
\import{\ProjectRootPath/preambles/custom}{list_setup}

% TODO: перевести команды на современный стиль объявления

\begin{document}

Для данных кривых будем находить дифференциальное уравнение, затем при выраженном $y\'\br{x} = f\br{x; y\br{x}}$ составим уравнение изогональных траекторий по правилу
$z\'\br{x} = \dfrac {f\br{x; z\br{x}} - 1} {1 + f\br{x; z\br{x}}}$ (см. задачи \textbf{40}, \textbf{46} и \textbf{48}). Полученные уравнения будут однородными.

\begin{enumerate}
    \item $y = x \ln\br{Cx}$.
    $C = \dfrac {\exp\br{\frac {y} {x}}} {x}$.
    $0 = \dfrac {\frac {y\'x - y} {\pow{x}{2}} \cdot \exp\br{\frac {y} {x}} \cdot x - \exp\br{\frac {y} {x}}} {\pow{x}{2}}$.
    $\dfrac {y\'x - y} {x} - 1 = 0$.
    $y\' = \dfrac {x + y} {x}$.
    $z\' = \dfrac {\dfrac {x + z} {x} - 1} {\dfrac {x + z} {x} + 1} = \dfrac {z} {2x + z}$.
    $u\br{x} = \dfrac {z\br{x}} {x}$.
    $z\' = u\'x + u = \dfrac {u} {2 + u}$.
    $u\'x = \dfrac {u - 2u - \pow{u}{2}} {u + 2}$.
    $\dfrac{u + 2}{\pow{u}{2} + u} u\' = -\dfrac {1} {x}$,
    причём $u\br{x} = 0$ и $u\br{x} = -1$ -- решения, \vspace{1mm} из которых получаются функции $z\br{x} = 0$ и $z\br{x} = -x$.

    $\dfrac {u + 2} {\pow{u}{2} + u} = \dfrac {2} {u} - \dfrac {1} {u + 1}$.
    $\tint \dfrac {u + 2} {\pow{u}{2} + u} u\' dx
    = \tint \br{\dfrac {2} {u} - \dfrac {1} {u + 1}} du
    = 2\ln\vbr{u} - \ln\vbr{u + 1} + C$. \vspace{1mm}

    $2\ln\vbr{u} - \ln\vbr{u + 1} = -\ln\vbr{x} + C$.
    $\pow{u}{2} \pow{x}{2} = C\br{ux + x}$.

    Итак, общая формула изогональных траекторий имеет вид
    $\pow{z}{2} = C\br{z + x}$,
    к которым необходимо добавить $z = -x$.

    \item $\br{x - 3y}^4 = Cx\pow{y}{6}$. \\
    $C = \dfrac {\br{x - 3y}^4} {x\pow{y}{6}}$.
    $0 = \dfrac {4\br{x - 3y}^3 \br{1 - 3y\'} \cdot x\pow{y}{6} - \br{x - 3y}^4 \br{\pow{y}{6} + 6x\pow{y}{5} y\'}} {\pow{x}{2} \pow{y}{12}}$. \\
    $4xy\br{1 - 3y\'} - \br{x - 3y}\br{y + 6xy\'} = 0$.
    $4xy - 12xyy\' - xy + 3\pow{y}{2} - 6\pow{x}{2}y\' + 18xyy\' = 0$.
    $xy + 2xyy\' + \pow{y}{2} - 2\pow{x}{2}y\' = 0$.
    $y\' = \dfrac {xy + \pow{y}{2}} {2\pow{x}{2} - 2xy}$.
    $z\' = \dfrac {-2\pow{x}{2} + 3xz + \pow{z}{2}} {2\pow{x}{2} - xz + \pow{z}{2}}$.
    $z\br{x} = x u\br{x}$. \vspace{1mm}
    $u\'x + u = \linebreak
    = \dfrac {-2 + 3u + \pow{u}{2}} {2 - u + \pow{u}{2}}$.
    $u\'x = \dfrac {\pow{u}{2} + 3u - 2 - 2u + \pow{u}{2} - \pow{u}{3}} {\pow{u}{2} - u + 2}
    = -\dfrac {\pow{u}{3} - 2\pow{u}{2} - u + 2} {\pow{u}{2} - u + 2}$. \vspace{2mm}
    Числитель разложим на множители методом группировки. \\
    $u\' x = -\dfrac {\pow{u}{2}\br{u - 2} - \br{u - 2}} {\pow{u}{2} - u + 2} = -\dfrac {\br{\pow{u}{2} - 1}\br{u - 2}}{\pow{u}{2} - u + 2} = -\dfrac {\br{u - 1}\br{u + 1}\br{u - 2}} {\pow{u}{2} - u + 2}$. \\
    $\dfrac {\pow{u}{2} - u + 2} {\br{u - 1}\br{u + 1}\br{u - 2}} u\' = - \dfrac {1} {x}$.
    Для взятия интеграла рациональную дробь в левой части разложим в сумму простейших методом вычёркивания.
    \begin{align*}
        \dfrac {\pow{u}{2} - u + 2} {\br{u - 1}\br{u + 1}\br{u - 2}} = \dfrac {a} {u - 1} + \dfrac{b} {u + 1} + \dfrac{c} {u - 2}\text{;}
        &&
        a = \dfrac {\pow{1}{2} - 1 + 2} {\br{1 + 1}\br{1 - 2}} = -1;
    \end{align*}
    \begin{align*}
        b = \dfrac {\br{-1}^2 - \br{-1} + 2} {\br{-1 - 1} \br{-1 - 2}} = \dfrac {2} {3};
        && c = \dfrac {\pow{2}{2} - 2 + 2} {\br{2 - 1} \br{2 + 1}} = \dfrac {4} {3};
    \end{align*}
    Итак,
    $\dfrac {\pow{u}{2} - u + 2} {\br{u - 1}\br{u + 1}\br{u - 2}} = - \dfrac {1} {u - 1} + \dfrac{2} {3\br{u + 1}} + \dfrac{4} {3\br{u - 2}}$. \\
    $\br{- \dfrac {1} {u - 1} + \dfrac{2} {3\br{u + 1}} + \dfrac{4} {3\br{u - 2}}} u\' = -\dfrac {1} {x}$.
    $\dfrac {2} {3} \ln\vbr{u + 1} + \dfrac {4} {3} \ln\vbr{u - 2} - \ln\vbr{u - 1} = -\ln\vbr{x} + C$.
    $\dfrac {\br{u + 1}^2\br{u - 2}^4} {\br{u - 1}^3} = \dfrac {C} {\pow{x}{3}}$.

    Получаем изогональные траектории в параметрической форме:
    $$\tsys{& x\br{u} = C f\br{u} \\& z\br{u} = C u f\br{u}} \text{, где } f\br{u} = \dfrac {u - 1} {\sqrt[3]{\br{u + 1}^2\br{u - 2}^4}}.$$

    Из соотношения
    $\dfrac {\br{u + 1}^2\br{u - 2}^4} {\br{u - 1}^3} = \dfrac {C} {\pow{x}{3}}$
    также можно получить уравнения изогональных траекторий в неявной форме $\br{u + 1}^2 \pow{x}{2} \cdot \br{u - 2}^4 \pow{x}{4} = C \br{u - 1}^3 \pow{x}{3}$.

    $\br{z + x}^2 \cdot \br{z - 2x}^4 = C \br{z - x}^3$.

    В процессе решения также были потеряны функции $u\br{x} = -1$, $u\br{x} = 2$ и $u\br{x} = 1$. При $C = 0$ первые две попадают в общую формулу, а последняя требует записи отдельно:
    $z\br{x} = x$.
\end{enumerate}

\end{document}

